\section{不完美信息下的算法细节：为什么 PPO 不够，Regret 家族更关键}

上一节建立了二人零和环境中 self-play 的理论地板；本节进一步追问：当玩家看不见完整状态、策略必须随机化时，什么算法才能真正逼近这块地板？阅读下面的公式与算法谱系时，应同时关注动作价值、混合概率和对手响应，不能把问题退化为普通单智能体策略优化。

\subsection{价值依赖混合概率，不只依赖动作本身}
Brown 通过 ``Rock-Paper-Scissors Plus'' 展示难点：某动作价值取决于其被选择概率。也就是说，算法不仅要学 ``做什么''，还要学 ``按什么频率做''。

形式化可写为：
\[
V(a_i) = \mathbb{E}_{a_{-i}\sim \pi_{-i}}[u(a_i, a_{-i})], \quad
\pi_i^\star = \arg\max_{\pi_i} \min_{\pi_{-i}} \mathbb{E}[u].
\]
在不完美信息博弈中，$\pi_i$ 的混合比例本身就是优化对象。

\begin{knowledgebox}{为什么 PPO 在这里会失效}
PPO 对单智能体 MDP 很强，但在需要精确概率混合与对手建模的场景，没有天然收敛到 Minimax 的保证，容易振荡或落入次优循环。
\end{knowledgebox}

\subsection{Fictitious Play 到 Regret Matching}
课程回顾了三类经典方法：
\begin{itemize}
    \item \textbf{Fictitious Play}：对历史平均策略做 best response，理论收敛但速度慢。
    \item \textbf{Regret Matching}：按正 regret 分配动作概率，显著加速。
    \item \textbf{Hedge/Regularized BR}：用正则化 best response 平衡收敛速度与稳定性。
\end{itemize}

\begin{importantbox}{算法演化主线}
从 ``精确 best response'' 走向 ``正则化 best response''，本质是用可控偏差换收敛速度和数值稳定性，这也是现代大规模博弈训练的核心工程策略。
\end{importantbox}

\subsection{扑克实战与方法选择}
前面的算法名称只有放回大规模扑克系统才显出差别：研究者面对的不是一个可枚举的小矩阵，而是由私有信息、公共历史和连续下注共同生成的巨大信息集。这里应比较的不是某次牌局胜负，而是算法能否在有限算力下持续压低可剥削性，并让搜索结果在真实对局中保持稳定。
Brown 提到其博士阶段扑克系统击败顶级职业选手，关键不是大模型，而是：
\begin{enumerate}
    \item 在不完美信息条件下做有效搜索；
    \item 在巨大策略空间内逼近低 exploitability 均衡。
\end{enumerate}
Regret 系方法是其中的关键部件。

\begin{warningbox}{工程告警：理论收敛不等于系统安全}
即使理论上可收敛，有限模型容量和有限训练步数也会留下 exploitable 缺口。实际部署前应做对抗探测与 exploitability 评估，而不是只看平均胜率。
\end{warningbox}

\subsection{本章小结}
不完美信息环境下，策略概率建模是第一公民。仅靠单智能体 RL 基线通常不足，Regret 家族方法更贴近问题结构。

